% Notas de aula — Capítulo 1: Fundamentos da Atmosfera
% Curso: Meteorologia Prática (baseado em R. Stull, Practical Meteorology, CC BY-NC-SA 4.0)
% Roteiro de quadro-negro:
%   - caixas verdes "Escrever no quadro" = conteúdo literal a escrever no quadro;
%   - linhas verdes com seta = instruções de encenação (desenhos, perguntas);
%   - caixas azul/vermelha = as duas formas de cada lei; caixa marrom = exemplos.
\documentclass[11pt]{article}
\usepackage{../comum/preambulo}
\graphicspath{{../figuras/cap01/}}

\title{Capítulo 1 --- Fundamentos da Atmosfera\\
{\large Notas de aula (roteiro de quadro-negro)}}
\author{Curso de Meteorologia Prática}
\date{}

\begin{document}
\maketitle
\thispagestyle{fancy}

\noindent\textbf{Plano sugerido:} 2 aulas de 2\,h.
\emph{Aula 1:} \S\ref{sec:intro}--\S\ref{sec:estrutura} (convenções, estado termodinâmico, estrutura vertical).
\emph{Aula 2:} \S\ref{sec:gas}--\S\ref{sec:processos} (lei dos gases, equilíbrio hidrostático, equação hipsométrica, processos).

\section{Introdução}\label{sec:intro}

\begin{noquadro}[abertura da aula]
\textbf{Cap.~1 --- Fundamentos da Atmosfera}\\[2pt]
Meteorologia $=$ física de um fluido (o ar) sobre uma esfera em rotação,
aquecida de forma desigual pelo Sol.\\[4pt]
Quatro ingredientes governam tudo:
\begin{enumerate}[itemsep=0pt,topsep=2pt]
  \item conservação de massa (continuidade);
  \item conservação de momento (2ª lei de Newton, referencial girante);
  \item conservação de energia (1ª lei da termodinâmica);
  \item equação de estado (lei dos gases ideais).
\end{enumerate}
\end{noquadro}

Vale gastar cinco minutos nesta lista: ela é o \emph{mapa do curso inteiro}.
A continuidade aparecerá quando falarmos de convergência e subsidência
\veremos{10}{divergência}; a 2ª lei de Newton num referencial girante é o
coração dos ventos \veremos{10}{força de Coriolis}; a 1ª lei da termodinâmica
governa o aquecimento e a convecção \veremos{3}{balanço de calor}; e a equação
de estado é usada já hoje. A mensagem: \emph{não há nada de exótico na
atmosfera} — é a física que os alunos já conhecem, aplicada a um sistema com
geometria e forçantes particulares.

Quase todo fenômeno meteorológico nasce de um único desequilíbrio: o Sol
aquece mais o equador que os polos, e a atmosfera (com os oceanos) transporta
esse excesso de energia \veremos{11}{circulação geral}. No caminho, o vapor
d'água condensa e libera calor latente \veremos{4}{umidade}
\veremos{7}{precipitação}, o ar gira por causa da rotação da Terra
\veremos{10}{ventos} e o resultado organizado chamamos de ``tempo''
\veremos{12}{massas de ar e frentes}.

A \textbf{composição do ar seco} (frações molares, homogênea até
$\sim$100\,km pela mistura turbulenta): N$_2$ 78{,}08\,\%, O$_2$ 20{,}95\,\%,
Ar 0{,}93\,\%, CO$_2$ $\sim$0{,}04\,\% (e crescendo $\sim$2,5\,ppm/ano — tema
do Cap.~21 \veremos{21}{mudança climática}). O vapor d'água varia de
$\sim 0$ a $4\,\%$: é o único constituinte que muda de fase nas temperaturas
terrestres, e por isso será tratado à parte \veremos{4}{vapor d'água}.

\quadro{Enfatizar: o livro do Stull usa formas algébricas (o público original
não tinha cálculo); nós escreveremos SEMPRE as duas formas lado a lado — a
algébrica do livro, para leitura e provas, e a diferencial, que conecta com a
mecânica dos fluidos que vocês conhecem.}

\section{Convenções meteorológicas}\label{sec:conv}

Antes de qualquer física, precisamos falar a língua da área. Dois pontos
pegam \emph{todo} iniciante: a direção do vento e o sistema de eixos.

\quadro{Desenhar uma rosa dos ventos com $0°=$N, $90°=$E, $180°=$S,
$270°=$W, e um vetor de vento entrando pela esquerda (vento de oeste).}

\begin{noquadro}[convenções que valem o curso inteiro]
\textbf{Direção do vento} $=$ direção \emph{de onde} o vento \emph{vem},
em graus, sentido horário a partir do Norte.\\
``Vento de oeste'' ($270°$) $\Rightarrow$ sopra de W para E.\\[4pt]
Eixos e componentes (sempre):\quad $x\to$ leste ($u$),\quad
$y\to$ norte ($v$),\quad $z\to$ cima ($w$).
\end{noquadro}

Por que a convenção ``de onde vem''? Herança náutica e prática: o vento que
\emph{vem} do mar traz umidade; o que vem do continente, ar seco — a origem
carrega a informação da massa de ar \veremos{12}{massas de ar}. Note o
contraste com a matemática: um ângulo medido no sentido horário a partir do
eixo $+y$, não anti-horário a partir de $+x$. A conversão entre as duas
convenções é fonte inesgotável de erros de código — o Caso~2 do notebook
existe para vacinar contra isso.

\begin{formalivro}[direção do vento a partir das componentes]
\[ \alpha = 90° - \frac{360°}{2\pi}\,\arctan\!\left(\frac{v}{u}\right) + \alpha_0,
\qquad \alpha_0 = 180° \text{ se } u>0,\; \alpha_0=0 \text{ se } u<0. \]
\end{formalivro}

\begin{formadif}[a mesma conta sem casos especiais]
Em código, usar sempre a arco-tangente de dois argumentos:
\[ \alpha = \operatorname{mod}\bigl(270° - \operatorname{atan2}(v,\,u)\cdot 180°/\pi,\ 360°\bigr), \]
que resolve os quadrantes automaticamente. Velocidade: $M=\sqrt{u^2+v^2}$.
\end{formadif}

\quadro{Exemplo rápido no quadro: $u=-5$\,m/s, $v=0$ $\Rightarrow$ o ar anda
para oeste, logo VEM do leste: vento de leste, $\alpha=90°$. Pedir mais dois
à turma ($u=0,v=5$; $u=v=3$).}

\section{Referenciais terrestres}\label{sec:terra}

Latitude $\phi$, longitude $\lambda$; raio médio da Terra
$R_T \simeq 6371$\,km. Um grau de latitude corresponde a
$2\pi R_T/360 \simeq 111$\,km — número que usaremos para converter gradientes
de pressão de cartas sinóticas em forças \veremos{10}{força do gradiente de
pressão}. Um grau de \emph{longitude} vale $111\cos\phi$\,km — encolhe até
zero nos polos.

\quadro{Desenhar observador, objeto no céu (sol, balão, satélite) e marcar os
três ângulos da caixa abaixo.}

\begin{noquadro}[ângulos locais]
Azimute $\alpha$: horizontal, horário a partir do Norte, até a projeção do
objeto no solo.\\
Elevação $\Psi$: acima do horizonte.\\
Zênite $\zeta$: a partir da vertical; $\ \zeta = 90° - \Psi$.
\end{noquadro}

Esses ângulos voltarão em dois papéis centrais: o ângulo zenital do Sol
comanda quanta radiação chega ao solo \veremos{2}{radiação} e as antenas de
radar apontam por azimute e elevação \veremos{8}{sensoriamento remoto}.

\textbf{Tempo}: toda observação meteorológica é registrada em UTC (tempo
universal coordenado) — uma rede mundial só funciona com um relógio só.
Cartas sinóticas são traçadas em 00, 06, 12 e 18\,UTC
\veremos{9}{cartas sinóticas}; o Brasil continental fica em UTC$-3$
(sem horário de verão desde 2019).

\section{Estado termodinâmico}\label{sec:estado}

\begin{noquadro}
O estado local do ar fica definido por três números: $(P,\ T,\ \rho)$\\
--- e uma equação de estado liga os três $\Rightarrow$ só dois são
independentes.
\end{noquadro}

Este é o mesmo enunciado da termodinâmica de equilíbrio: um gás simples tem
dois graus de liberdade termodinâmicos. A umidade acrescentará uma quarta
variável \veremos{4}{razão de mistura}, e a meteorologia inventará dezenas de
``temperaturas'' derivadas ($T_v$ hoje; potencial $\theta$ no Cap.~3; de
bulbo úmido, equivalente\ldots\ no Cap.~4) — todas são apenas combinações
convenientes dessas variáveis básicas.

\subsection{Temperatura}

Temperatura mede o movimento molecular: energia cinética média de
$\tfrac12 k_B T$ por grau de liberdade. A rigor, portanto, $T$ só está
definida para um \emph{conjunto} grande de moléculas — o que em 1\,mm$^3$ de
ar ($\sim 10^{16}$ moléculas) nunca é problema.

\begin{formalivro}[conversões]
\[ T\,[\mathrm{K}] = T\,[°\mathrm{C}] + 273{,}15, \qquad
   T\,[°\mathrm{C}] = \tfrac{5}{9}\,(T\,[°\mathrm{F}] - 32). \]
\end{formalivro}

\quadro{Avisar com ênfase: TODA fórmula física do curso usa kelvin. O erro
de usar °C numa lei de gases ou numa razão de temperaturas é o erro numérico
mais comum do semestre.}

Valores para calibrar a intuição: padrão ao nível do mar
$T=15\,°$C$\,=288{,}15$\,K; recorde de calor $\sim 57\,°$C; de frio
$\sim -89\,°$C; tropopausa tropical $\sim -80\,°$C
\javimos{1}{atmosfera padrão, \S\ref{sec:estrutura}}.

\subsection{Pressão}

Pressão é força normal por área, das colisões moleculares. Ponto sutil que o
Stull destaca (Fig.~1.7 do livro): num fluido em repouso a pressão é
\emph{isotrópica} — empurra igualmente em todas as direções, inclusive para
cima. É por isso que a pressão pode sustentar o peso do ar de cima
(\S\ref{sec:hidro}) e por que não sentimos os $\sim 10^4$\,kg de ar sobre
cada m$^2$ dos nossos ombros.

\begin{formalivro}[unidades]
\[ 1\ \mathrm{kPa} = 10\ \mathrm{hPa} = 10\ \mathrm{mb};\qquad
P_{\text{nível do mar}} = 101{,}325\ \mathrm{kPa} \simeq 1013{,}25\ \mathrm{hPa}. \]
\end{formalivro}

O livro usa kPa; a prática operacional (METAR, cartas) usa hPa
\veremos{9}{METAR}. Cobrar fluência nas duas — o fator é só $10\times$, mas
em prova ele derruba gente.

A pressão decai quase exponencialmente com a altura:

\begin{formalivro}[perfil de pressão, Stull eq.~1.9]
\[ P \simeq P_0\, e^{-z/H_p}, \qquad H_p = 7{,}29\ \mathrm{km}
\ \ (\text{escala de altura da pressão}). \]
\end{formalivro}

\begin{formadif}[de onde vem a exponencial]
Combinando equilíbrio hidrostático ($dP/dz=-\rho g$, \S\ref{sec:hidro}) com a
lei dos gases ($\rho = P/\Re_d T$, \S\ref{sec:gas}) numa atmosfera isotérmica:
\[ \frac{dP}{P} = -\frac{g}{\Re_d T}\,dz
\quad\Longrightarrow\quad
P(z) = P_0\, \exp\!\left(-\frac{z}{H}\right),\qquad H = \frac{\Re_d T}{g}. \]
Com $T\simeq250$\,K (média da troposfera), $H\simeq7{,}3$\,km — o valor
empírico do livro. A atmosfera real não é isotérmica, por isso a exponencial
é só uma (excelente) aproximação.
\end{formadif}

\begin{noquadro}[números para levar de cor]
$P$ cai à metade a cada $\sim 5{,}5$\,km\quad
($H_p \ln 2 = 7{,}29\times0{,}693$).\\
Em $z=5{,}5$\,km: $\sim 50$\,kPa (metade da atmosfera abaixo de você).\\
Everest ($8{,}8$\,km): $\sim 30$\,kPa. Cruzeiro de avião ($11$\,km):
$\sim 22$\,kPa.
\end{noquadro}

\quadro{Desenhar $P(z)$ em escala linear e, ao lado, em escala mono-log — a
reta no gráfico log é o teste visual da exponencial. (Caso~1 do notebook
refaz este gráfico com a atmosfera padrão ICAO.)}

\subsection{Densidade}

\begin{formalivro}[perfil de densidade, Stull eq.~1.13]
\[ \rho \simeq \rho_0\, e^{-z/H_\rho}, \qquad
\rho_0 = 1{,}2250\ \mathrm{kg/m^3},\quad H_\rho = 8{,}55\ \mathrm{km}. \]
\end{formalivro}

Pergunta para a turma: por que $H_\rho \neq H_p$? (Resposta: $\rho=P/\Re_d T$;
como $T$ também cai com a altura na troposfera, $\rho$ cai \emph{mais devagar}
que $P$. Se $T$ fosse constante, as duas escalas coincidiriam — é exatamente
o exercício T1.)

A densidade é a variável ``invisível'' — não há densímetro de ar nas
estações; ela é sempre \emph{calculada} da lei dos gases. Mas é ela que entra
no empuxo \veremos{5}{estabilidade e flutuabilidade} e na potência extraível
do vento \veremos{18}{energia eólica}.

\section{Estrutura vertical da atmosfera}\label{sec:estrutura}

\subsection{Atmosfera padrão}

\quadro{Desenhar o perfil $T(z)$ da atmosfera padrão de 0 a 50\,km, com as
camadas e as pausas nomeadas. Este desenho fica no quadro até o fim da aula —
tudo que vier depois se pendura nele.}

\figlivro{fig_01_10.jpg}{0.5}{Atmosfera padrão: perfil vertical de
temperatura com as camadas atmosféricas.}{Fig.~1.10}

A ``atmosfera padrão'' ICAO é uma atmosfera de referência, por convenção
internacional: útil para calibrar altímetros, projetar aeronaves e — para
nós — como estado de fundo contra o qual comparar sondagens reais
\veremos{5}{sondagens e diagramas termodinâmicos}.

\begin{noquadro}[atmosfera padrão ICAO]
\begin{tabular}{lll}
Camada & Intervalo & Gradiente \\
\hline
Troposfera & 0--11\,km & $-6{,}5$\,K/km (de $288{,}15$\,K) \\
Tropopausa/estrat.\ inferior & 11--20\,km & isotérmica, $216{,}65$\,K \\
Estratosfera média & 20--32\,km & $+1{,}0$\,K/km \\
Estratosfera superior & 32--47\,km & $+2{,}8$\,K/km \\
Estratopausa & 47--51\,km & isotérmica, $270{,}65$\,K \\
\end{tabular}
\end{noquadro}

\begin{formalivro}[troposfera padrão]
\[ T(z) = T_0 - \gamma\, z, \qquad T_0 = 288{,}15\ \mathrm{K},\quad
\gamma = 6{,}5\ \mathrm{K/km}\ (0 \le z \le 11\ \mathrm{km}). \]
\end{formalivro}

Atenção ao sinal: $\gamma$ (\emph{lapse rate}, taxa de decaimento) é definido
\emph{positivo} quando $T$ cai com a altura. Guardar este $6{,}5$\,K/km: ele
será comparado com os gradientes adiabáticos seco ($9{,}8$\,K/km) e úmido
($\sim 4$--$7$\,K/km) para decidir estabilidade \veremos{3}{gradiente
adiabático} \veremos{5}{estabilidade estática}.

\subsection{Camadas e suas físicas}

O perfil de $T$ não é decorativo: cada mudança de sinal do gradiente denuncia
\emph{onde a atmosfera absorve energia}.

\begin{itemize}
  \item \textbf{Troposfera} (0--11\,km): contém $\sim 80\%$ da massa e
  praticamente todo o ``tempo''. $T$ cai com $z$ porque a fonte de calor é a
  \emph{superfície} (aquecida pelo Sol) \veremos{2}{balanço radiativo};
  aquecida por baixo, ela convecta como uma panela no fogo
  \veremos{3}{convecção}.
  \item \textbf{Estratosfera} (11--51\,km): $T$ \emph{cresce} com $z$ porque o
  ozônio absorve UV — uma fonte de calor \emph{no topo}. Aquecida por cima,
  é estratificada e estável (daí o nome): convecção suprimida, nuvens raras,
  é onde voam os jatos comerciais fugindo da turbulência
  \veremos{5}{estabilidade}.
  \item \textbf{Mesosfera} (51--85\,km): sem absorvedor, $T$ volta a cair.
  \textbf{Termosfera} (acima): UV extremo, $T$ sobe de novo (mas com
  densidade tão baixa que ``temperatura'' ali quase não aquece nada).
  \item \textbf{Camada limite atmosférica} (0--1\,km típico): a porção da
  troposfera que responde à superfície em $\lesssim 1$\,h, com ciclo diurno
  forte. É onde vivemos; ganha um capítulo próprio \veremos{18}{camada
  limite}.
\end{itemize}

\subsection{Um comentário sobre pesos e medidas}

A coluna inteira de ar pesa $P_0/g \simeq 10^4$\,kg por m$^2$ — dez
toneladas sobre cada metro quadrado (exercício T4 pede a massa total da
atmosfera). A metade de baixo cabe nos primeiros $5{,}5$\,km: a atmosfera é
uma película — espessura de $\sim 0{,}1\%$ do raio terrestre — e é essa
razão de aspecto extrema que tornará excelentes as aproximações hidrostática
(\S\ref{sec:hidro}) e ``rasa'' da dinâmica \veremos{10}{aproximações da
dinâmica}.

\section{Equação de estado --- lei dos gases ideais}\label{sec:gas}

\begin{noquadro}[a equação mais usada do curso]
\[ \boxed{\,P = \rho\, \Re_d\, T\,}\qquad
\Re_d = 287\ \mathrm{J\,kg^{-1}\,K^{-1}} =
0{,}287\ \mathrm{kPa\,K^{-1}\,m^3\,kg^{-1}} \]
vale para ar seco; para ar úmido, trocar $T \to T_v$ (abaixo).
\end{noquadro}

\begin{formadif}[conexão com a física estatística]
Da forma molecular $PV = N k_B T$, dividindo pela massa $M = N\bar m$:
\[ P = \frac{N k_B T}{V} = \rho\,\frac{k_B}{\bar m}\,T \equiv \rho\,\Re\,T, \]
ou seja, $\Re = k_B/\bar m = R^*/M_{\text{molar}}$: constante universal
$R^* = 8{,}314$\,J\,mol$^{-1}$K$^{-1}$ dividida pela massa molar média do ar
seco ($28{,}97$\,g/mol) $\Rightarrow \Re_d = 287$\,J\,kg$^{-1}$K$^{-1}$.
A meteorologia prefere a forma com $\rho$ (em vez de $V$) porque o ar não tem
``recipiente'': trabalhamos com campos locais e intensivos.
\end{formadif}

Por que podemos tratar o ar como ideal? As correções de gás real (van der
Waals) escalam com a razão entre o volume das moléculas e o volume
disponível; nas condições atmosféricas isso dá desvios $<0{,}2\%$ — menor
que a incerteza dos termômetros operacionais.

\subsection{Temperatura virtual}

Fato contraintuitivo que decide tempestades: ar úmido é \emph{menos} denso
que ar seco à mesma $(P,T)$, porque a molécula de água (18\,g/mol) é mais
leve que a média do ar (29\,g/mol) — cada H$_2$O que entra na mistura
\emph{substitui} um N$_2$ ou O$_2$. Em vez de trocar a constante $\Re$ a cada
umidade, a meteorologia corrige a temperatura:

\begin{formalivro}[Stull eq.~1.21]
\[ T_v = T\,\bigl(1 + 0{,}61\, r - r_L - r_I\bigr), \]
com $r$ = razão de mistura de vapor, $r_L,\ r_I$ = razões de mistura de água
líquida e gelo (g/g) \veremos{4}{razão de mistura}. Então
$P=\rho\,\Re_d\,T_v$ vale para ar úmido e até nublado.
\end{formalivro}

Leitura da fórmula: o vapor ($+0{,}61r$) \emph{alivia} o ar; gotículas e
cristais em suspensão ($-r_L-r_I$) são lastro que o \emph{pesa}. $T_v$ é a
temperatura que o ar \emph{seco} precisaria ter para ficar com a densidade do
ar úmido real.

\quadro{Comentar: a correção é de no máximo 2--3\,K, mas é exatamente essa
diferença que decide se uma parcela flutua ou afunda — o empuxo será definido
com $T_v$ no Cap.~5.}\veremos{5}{CAPE e empuxo}

\begin{exemplo}[temperatura virtual]
Ar a $T=25\,°$C com $r=20$\,g/kg $=0{,}020$\,g/g, sem condensados.
\[ T_v = (298{,}15\ \mathrm{K})\times(1+0{,}61\times0{,}020)
       = 298{,}15\times1{,}0122 = 301{,}8\ \mathrm{K} = 28{,}6\,°\mathrm{C}. \]
Para a dinâmica, esse ar úmido se comporta como ar seco $3{,}6$\,K mais
quente. \emph{Verificação: $r$ adimensional em g/g; correção de poucos K,
magnitude típica OK.}
\end{exemplo}

\section{Equilíbrio hidrostático}\label{sec:hidro}

\quadro{Desenhar uma ``fatia'' de ar de área $A$ e espessura $\Delta z$, com
as três forças: $P(z)A$ para cima, $P(z+\Delta z)A$ para baixo, peso
$\rho A\Delta z\,g$ para baixo. A dedução inteira sai deste desenho.}

Por que o ar não cai? Porque a pressão embaixo é maior que em cima
(justamente porque há mais ar acima comprimindo), e a diferença sustenta o
peso. Escrever esse balanço é a dedução mais importante da aula:

\begin{formadif}[dedução no quadro]
\[ \underbrace{P(z)A}_{\text{empurra p/ cima}}
 - \underbrace{P(z+\Delta z)A}_{\text{empurra p/ baixo}}
 - \underbrace{\rho A\,\Delta z\, g}_{\text{peso}} = 0
\quad\Longrightarrow\quad
\boxed{\ \frac{dP}{dz} = -\rho\, g\ } \]
com $g=9{,}8$\,m\,s$^{-2}$.
\end{formadif}

\begin{formalivro}[Stull eq.~1.25]
\[ \frac{\Delta P}{\Delta z} \simeq -\rho\,g. \]
Regra prática perto da superfície: $P$ cai $\sim1$\,kPa a cada $\sim84$\,m
de subida ($\rho\simeq1{,}2$\,kg/m$^3$).
\end{formalivro}

\begin{noquadro}[para fixar]
hidrostática $\equiv$ ``a pressão num nível é o peso, por área, de todo o ar
acima'':\quad $P(z) = g\displaystyle\int_z^\infty \rho\,dz'$.
\end{noquadro}

\begin{exemplo}[queda de pressão por andar]
Prédio com $3$\,m por andar, $\rho=1{,}2$\,kg/m$^3$:
\[ \Delta P = -\rho g\,\Delta z = -(1{,}2)(9{,}8)(3) \simeq -35\ \mathrm{Pa}
 = -0{,}035\ \mathrm{kPa}. \]
O barômetro MEMS de um celular resolve $\sim10$\,Pa: dá para ``pesar'' um
andar em sala de aula — o exercício N3 usa dados reais de elevador.
\end{exemplo}

\emph{Limite de validade}: a hidrostática assume aceleração vertical
desprezível. Vale para tudo que é largo e raso (escala sinótica: 1000\,km de
largura por 10\,km de altura); falha \emph{dentro} de correntes convectivas
intensas, onde $w$ muda rápido \veremos{14}{tempestades}. Os modelos
numéricos globais são hidrostáticos; os de tempestade, não
\veremos{20}{previsão numérica}.

\section{Equação hipsométrica}\label{sec:hipso}

\quadro{Anunciar o objetivo: ``medir altura com um barômetro''. As duas
equações-mestras já estão no quadro (hidrostática + gases) — a hipsométrica
é o primeiro fruto do casamento delas.}

\begin{formadif}[dedução]
\[ \frac{dP}{dz} = -\rho g, \qquad \rho = \frac{P}{\Re_d T_v}
\quad\Longrightarrow\quad
\frac{dP}{P} = -\frac{g}{\Re_d T_v}\,dz. \]
Integrando de $(z_1,P_1)$ a $(z_2,P_2)$ com $\bar T_v$ = média da camada:
\[ \boxed{\ z_2 - z_1 = \frac{\Re_d\,\bar T_v}{g}\,
   \ln\!\left(\frac{P_1}{P_2}\right)\ } \]
\end{formadif}

\begin{formalivro}[Stull eq.~1.26]
\[ z_2 - z_1 = a\,\bar T_v \ln\!\left(\frac{P_1}{P_2}\right), \qquad
   a = \frac{\Re_d}{g} = 29{,}3\ \mathrm{m/K}. \]
\end{formalivro}

\begin{noquadro}[interpretação --- escrever por extenso]
A espessura de uma camada entre duas isóbaras é proporcional à sua
temperatura média:\\
\textbf{camada quente $=$ camada espessa; camada fria $=$ camada rasa.}
\end{noquadro}

Esta frase simples é uma das mais férteis do curso: dela sairão a carta de
espessura como termômetro operacional \veremos{9}{análise de cartas}, a
inclinação das superfícies de pressão com a latitude e, com Coriolis, o vento
térmico e a existência do jato \veremos{11}{vento térmico e jato}. Também é o
princípio do altímetro barométrico de qualquer aeronave.

\begin{exemplo}[espessura da camada 100--50 kPa]
$\bar T_v = 260$\,K entre $P_1=100$\,kPa e $P_2=50$\,kPa:
\[ \Delta z = (29{,}3\ \mathrm{m/K})(260\ \mathrm{K})\ln 2
 = 29{,}3\times260\times0{,}693 \simeq 5280\ \mathrm{m}. \]
A ``espessura 1000--500'' é usada como termômetro da troposfera inferior;
regra clássica: $5400$\,m separa chuva de neve nas latitudes médias.
\emph{Verificação: $\ln 2=0{,}693$; resultado coerente com
$H_p\ln 2\simeq5$\,km.}
\end{exemplo}

\section{Terminologia de processos}\label{sec:processos}

\quadro{Construir a tabela no quadro de dois em dois, pedindo exemplos à
turma antes de escrever a coluna da direita.}

\begin{noquadro}[dicionário de processos]
\begin{tabular}{ll}
isobárico & $P$ constante \\
isotérmico & $T$ constante \\
isostérico & volume específico $1/\rho$ constante \\
adiabático & sem troca de calor \\
diabático & com troca de calor \\
isentrópico & entropia constante (adiabático reversível) \\
\end{tabular}
\end{noquadro}

Conceito central para o próximo capítulo: a subida de uma parcela de ar é,
em excelente aproximação, \emph{adiabática} — o ar é péssimo condutor
térmico e a mistura com o ambiente é lenta comparada ao tempo de subida
\veremos{3}{processo adiabático e $\theta$}. ``Diabático'' então é o resto:
radiação \veremos{2}{aquecimento radiativo} e calor latente
\veremos{4}{calor latente} \veremos{7}{precipitação}.

\section{Instrumentos de pressão}

Barômetro de mercúrio (referência histórica: $760$\,mm\,Hg
$=101{,}325$\,kPa; Torricelli, 1643), aneróide (cápsula que deforma) e
sensores eletrônicos capacitivos/MEMS, hoje em qualquer celular e estação
automática. As estações reportam a pressão \emph{reduzida ao nível do mar}
via equação hipsométrica \javimos{1}{eq.\ hipsométrica, \S\ref{sec:hipso}} —
sem essa redução, um mapa de pressão seria só um mapa de altitude das
estações \veremos{9}{redução ao nível do mar}.

\section{Síntese da aula}

\begin{noquadro}[mapa final --- deixar num canto do quadro]
\[ P = \rho\,\Re_d T_v \qquad \frac{dP}{dz} = -\rho g \]
\[ z_2 - z_1 = a\,\bar T_v\ln\frac{P_1}{P_2} \qquad P \simeq P_0\,e^{-z/H_p} \]
duas leis (estado, hidrostática) $+$ duas consequências (hipsométrica,
exponencial).
\end{noquadro}

No notebook do capítulo os alunos verificam numericamente as quatro contra a
atmosfera padrão ICAO (pacote \texttt{ambiance}) e contra uma sondagem real
de Campo de Marte (pacote \texttt{MetPy}) \veremos{5}{sondagens}.

\section{Exercícios complementares}\label{sec:exercicios}

Complementam os exercícios do livro (seção 1.12). Os teóricos (T) são para
entrega escrita; os numéricos (N) resolvem-se no notebook
\texttt{cap01\_fundamentos.ipynb}. Soluções completas (para o professor) em
\texttt{cap01\_solucoes.ipynb}.

\subsection*{Teóricos}

\begin{enumerate}[label=\textbf{T\arabic*.},itemsep=4pt]
  \item Mostre que numa atmosfera \emph{isotérmica} as escalas de altura de
        pressão e densidade coincidem, $H_p = H_\rho = \Re_d T/g$, e
        calcule-as para $T=250$\,K. Explique por que na atmosfera real
        $H_\rho > H_p$.
  \item Para uma troposfera com $T(z)=T_0-\gamma z$, deduza o perfil
        politrópico
        $P(z)=P_0\,(1-\gamma z/T_0)^{\,g/(\Re_d\gamma)}$
        e mostre que no limite $\gamma\to0$ ele recupera a exponencial
        isotérmica.
  \item A partir da lei de Dalton para a mistura ar seco + vapor, com
        $\varepsilon = M_v/M_d = 0{,}622$, deduza
        $T_v \simeq T\,(1+0{,}61\,r)$ em primeira ordem na razão de mistura
        $r$.
  \item Mostre, integrando a equação hidrostática, que a massa de ar por
        unidade de área da superfície é $m = P_0/g$, e estime a massa total
        da atmosfera terrestre.
  \item Mostre que a espessura da camada 1000--500\,hPa vale
        $\Delta z = a\,\bar T_v \ln 2$ e determine o $\bar T_v$
        correspondente ao limiar operacional de 5400\,m. Discuta o sentido
        físico da regra chuva/neve.
\end{enumerate}

\subsection*{Numéricos (no notebook)}

\begin{enumerate}[label=\textbf{N\arabic*.},itemsep=4pt]
  \item Integre numericamente $dP/dz = -gP/(\Re_d T)$ com o perfil
        politrópico do T2 (0 a 11\,km) e compare com a solução analítica;
        trace o erro relativo.
  \item Ajuste por mínimos quadrados a melhor escala de altura $H_p$ da
        atmosfera padrão (0--20\,km) e ache as altitudes em que $P$ vale
        $\tfrac12$, $\tfrac14$ e $\tfrac1{10}$ de $P_0$.
  \item Com a série temporal de pressão de um barômetro dentro de um
        elevador (dados no notebook), reconstrua $z(t)$ pela equação
        hipsométrica; determine a altura do prédio e a velocidade média.
  \item Para a sondagem do Caso~4, calcule a espessura 1000--500\,hPa com o
        MetPy e aplique a regra dos 5400\,m; repita esfriando a sondagem em
        30\,K.
\end{enumerate}

\vspace{1em}
\noindent\rule{\linewidth}{0.4pt}\\
{\scriptsize Material derivado de R.~Stull, \emph{Practical Meteorology}
(CC BY-NC-SA 4.0); este documento herda a mesma licença.}

\end{document}
